
\documentclass[12pt,a4paper]{report}
\usepackage{fontspec}
\setmainfont{Times New Roman}
\setmonofont{Courier New}
\usepackage[top=2.5cm,bottom=2.5cm,left=3cm,right=2cm]{geometry}
\usepackage{amsmath,amssymb}
\usepackage{graphicx}
\usepackage{booktabs,tabularx,array,multirow}
\usepackage{float}
\usepackage{xcolor}
\usepackage{tikz}
\usetikzlibrary{arrows.meta,positioning,shapes.geometric}
\usepackage{hyperref}
\usepackage{fancyhdr}
\usepackage{titlesec}
\usepackage{caption}
\usepackage{enumitem}
\usepackage{fvextra}
\usepackage{needspace}

\definecolor{reportblue}{RGB}{0,45,145}
\definecolor{softblue}{RGB}{237,244,251}
\hypersetup{colorlinks=true,linkcolor=reportblue,citecolor=reportblue,
  urlcolor=reportblue,pdftitle={Agent 02: RL agent va rule-based agent},
  pdfauthor={Phan Van Khoi}}
\renewcommand{\chaptername}{Phần}
\renewcommand{\contentsname}{Mục lục}
\renewcommand{\listfigurename}{Danh sách hình}
\renewcommand{\listtablename}{Danh sách bảng}
\renewcommand{\figurename}{Hình}
\renewcommand{\tablename}{Bảng}
\renewcommand{\bibname}{Tài liệu tham khảo}
\titleformat{\chapter}[display]
  {\normalfont\bfseries\color{reportblue}}
  {\filleft\Large\chaptername\ \thechapter}{1ex}
  {\titlerule\vspace{1ex}\filright\LARGE}
  [\vspace{1ex}\titlerule]
\titleformat{\section}{\normalfont\large\bfseries\color{reportblue}}{\thesection}{0.7em}{}
\titleformat{\subsection}{\normalfont\normalsize\bfseries\color{reportblue}}{\thesubsection}{0.6em}{}
\setlength{\parindent}{1.1cm}
\setlength{\parskip}{0.28em}
\linespread{1.16}
\setlist[itemize]{leftmargin=1.8em,itemsep=0.18em,topsep=0.25em}
\setlist[enumerate]{leftmargin=1.8em,itemsep=0.18em,topsep=0.25em}
\captionsetup{font=small,labelfont=bf}
\DefineVerbatimEnvironment{Code}{Verbatim}{fontsize=\scriptsize,frame=single,
  breaklines=true,breakanywhere=true,baselinestretch=1.0}
\DefineVerbatimEnvironment{Output}{Verbatim}{fontsize=\scriptsize,frame=single,
  breaklines=true,breakanywhere=true}
% Code extracts are read directly from the source file, so the report
% always shows the code that produced the results.
\newcommand{\SRC}{"../../src/Agent 02/agent02_experiment.py"}
\newcommand{\CodeLines}[2]{\VerbatimInput[fontsize=\scriptsize,frame=single,
  breaklines=true,numbers=left,numbersep=4pt,firstline=#1,lastline=#2,
  baselinestretch=1.0]{\SRC}}
\newcommand{\Fig}[3][0.92]{\begin{figure}[H]\centering
  \includegraphics[width=#1\textwidth]{images/#2}\caption{#3}\end{figure}}
\pagestyle{fancy}
\fancyhf{}
\fancyhead[R]{\small\itshape Agent 02: RL agent và rule-based agent}
\fancyfoot[C]{\thepage}
\setlength{\headheight}{15pt}
\setlength{\emergencystretch}{3em}

\begin{document}
\hypersetup{pageanchor=false}
\begin{titlepage}
\thispagestyle{empty}
\begin{center}
\vspace*{0.7cm}
{\bfseries\fontsize{16}{20}\selectfont HỌC VIỆN CÔNG NGHỆ BƯU CHÍNH VIỄN THÔNG}\\[0.15cm]
{\bfseries\fontsize{15}{19}\selectfont KHOA CÔNG NGHỆ THÔNG TIN 1}\\[0.3cm]
\rule{5cm}{0.5pt}
\vspace{2.2cm}

{\bfseries\fontsize{18}{24}\selectfont BÁO CÁO THÍ NGHIỆM AGENT 02}\\[0.35cm]
{\bfseries\fontsize{17}{23}\selectfont PHÁT TRIỂN CÁC HỆ THỐNG THÔNG MINH}\\[0.7cm]
{\bfseries\color{reportblue}\fontsize{21}{27}\selectfont AGENT HỌC TĂNG CƯỜNG\\[0.15cm]
VÀ AGENT DỰA TRÊN LUẬT}\\[0.3cm]
{\bfseries\fontsize{15}{20}\selectfont Mô phỏng hai người bán cạnh tranh giảm giá trên sàn thương mại điện tử}

\vspace{2.0cm}
\begin{tabular}{>{\bfseries}p{4.1cm}p{6.5cm}}
Họ tên: & Phan Văn Khôi \\
Mã sinh viên: & B23DCCN447 \\
Lớp: & D23CTPM01-B \\
Giảng viên: & PGS. TS. Trần Đình Quế \\
\end{tabular}
\vfill
{\bfseries\fontsize{14}{18}\selectfont Hà Nội -- 2026}
\end{center}
\end{titlepage}
\hypersetup{pageanchor=true}
\pagenumbering{roman}
\tableofcontents
\clearpage
\listoftables
\listoffigures
\clearpage
\pagenumbering{arabic}

% ============================================================
\chapter{Giới thiệu và mục tiêu}
\label{chap:intro}

Thí nghiệm Agent 02 so sánh hai cách xây dựng agent ra quyết định trong cùng một
môi trường: \textbf{agent học tăng cường} (RL) dùng Q-learning dạng bảng, và
\textbf{agent dựa trên luật} (rule-based) dùng chính sách if--else do người viết sẵn
\cite{assignment,complementary}. Hai người bán cùng bán một loại sản phẩm; ở mỗi bước
thời gian, mỗi người quyết định \emph{giảm giá} hoặc \emph{không giảm giá}. Giảm giá
làm tăng xác suất khách mua nhưng làm giảm lợi nhuận trên mỗi sản phẩm, nên agent
phải cân bằng hai tác động này.

Mục tiêu của báo cáo:
\begin{enumerate}
\item Mô hình hóa bài toán theo các khái niệm agent, môi trường, trạng thái, hành
  động, phần thưởng và chính sách; giải thích vì sao đây là hệ nhiều agent.
\item Cài đặt môi trường, agent Q-learning và agent dựa trên luật bằng Python, tách
  riêng giai đoạn huấn luyện và giai đoạn đánh giá.
\item Thực hiện năm thí nghiệm A--E của đề bài: chạy baseline, xem quá trình học,
  thay đổi tốc độ học, thay đổi mức khám phá và so sánh công bằng trên nhiều seed.
\item Trả lời mười câu hỏi bằng số liệu đo được; cài đặt và đánh giá một thay đổi
  môi trường (câu 9) và một agent cải tiến (câu 10).
\end{enumerate}

\noindent\textbf{Mã nguồn trên GitHub:} \url{https://github.com/KhoiF/PTHTTM}.
Chương trình nằm tại \path{src/Agent 02/agent02_experiment.py}; toàn bộ kết quả
(CSV, PNG, bảng Q và nhật ký chạy \path{run_log.txt}) nằm trong
\path{src/Agent 02/agent02_results/}. Mọi con số trong báo cáo được lấy từ các tệp
này, không có số liệu ước đoán. Các đoạn code trong báo cáo được trích tự động từ
tệp nguồn khi biên dịch LaTeX, số dòng ở lề trái là số dòng trong tệp.

\paragraph{Kết luận chính.} Trong môi trường tham chiếu (giảm giá 20\%), chính sách
tối ưu khi đối thủ là agent dựa trên luật là \emph{không bao giờ giảm giá}, đạt
lợi nhuận kỳ vọng 35{,}00 mỗi bước. Agent dựa trên luật cũng không bao giờ giảm giá
trong tình huống này nên đã đạt đúng mức tối ưu. Agent RL học được chính sách
tối ưu ở 9/10 seed, đạt trung bình $34{,}72\pm0{,}75$; chênh lệch so với agent dựa
trên luật không có ý nghĩa thống kê. Khi đổi mức giảm giá xuống 5\%, luật cố định
không còn tối ưu và RL vượt luật ($36{,}72$ so với $34{,}99$). Agent dựa trên mô hình
đề xuất ở câu 10 gần đạt tối ưu trong cả sáu kịch bản thử.

% ============================================================
\chapter{Môi trường và hai agent}
\label{chap:model}

\section{Trạng thái, hành động và động lực môi trường}
Với agent $i\in\{1,2\}$, trạng thái cục bộ tại thời điểm $t$ là
\begin{equation}
s_i(t)=(d_t,\;c_i(t)),\qquad d_t\in\{0,1\},\; c_i(t)=a_{-i}(t-1)\in\{0,1\},
\end{equation}
trong đó $d_t$ là nhu cầu thị trường (0: thấp, 1: cao) và $c_i(t)$ là hành động
của đối thủ ở bước trước. Mỗi agent có 4 trạng thái và 2 hành động
$a_i(t)\in\{0,1\}$ (0: không giảm giá, 1: giảm giá). Hai agent chọn hành động
\emph{đồng thời}, nên không agent nào biết hành động hiện tại của đối thủ.

Xác suất một khách tiềm năng mua hàng của agent $i$ là
\begin{equation}
p_i(t)=\operatorname{clip}\bigl(0{,}30+0{,}25d_t+0{,}15a_i(t)-0{,}10a_{-i}(t),\;0{,}05,\;0{,}95\bigr).
\label{eq:conv}
\end{equation}
Số khách tiềm năng là $N_t=10$ khi $d_t=0$ và $N_t=20$ khi $d_t=1$. Số đơn bán được
$S_i(t)\sim\operatorname{Binomial}(N_t,p_i(t))$. Với giá gốc $P=10$, giá vốn $C=5$
và mức giảm $\delta=0{,}20$, phần thưởng là lợi nhuận
\begin{equation}
r_i(t)=S_i(t)\bigl[P(1-\delta a_i(t))-C\bigr].
\label{eq:reward}
\end{equation}
Sau khi trả thưởng, môi trường sinh $d_{t+1}$ ngẫu nhiên đều, độc lập với hành động,
và hành động hiện tại trở thành ``hành động trước'' của bước sau. Mỗi episode có
20 bước; đầu episode, nhu cầu được rút ngẫu nhiên và hành động trước của cả hai
agent được đặt bằng 0.

\begin{figure}[H]
\centering
\begin{tikzpicture}[node distance=6mm and 10mm,>=Stealth,
  box/.style={draw=reportblue,rounded corners,fill=softblue,
  align=center,minimum height=10mm,text width=3.3cm,font=\small}]
\node[box] (env) {Môi trường\\ $d_t$, $(a_1,a_2)_{t-1}$};
\node[box,above left=8mm and 4mm of env] (a1) {Agent 1 (RL)\\ quan sát $(d_t,a_2(t-1))$};
\node[box,above right=8mm and 4mm of env] (a2) {Agent 2 (luật)\\ quan sát $(d_t,a_1(t-1))$};
\node[box,below=of env,text width=6.5cm] (out) {Doanh số $S_i\sim\mathrm{Bin}(N_t,p_i)$, thưởng $r_i$;\\ sinh $d_{t+1}$ độc lập với hành động};
\draw[->,thick] (env) -- node[left,font=\scriptsize]{$s_1(t)$} (a1);
\draw[->,thick] (env) -- node[right,font=\scriptsize]{$s_2(t)$} (a2);
\draw[->,thick] (a1.south) to[bend right=25] node[left,font=\scriptsize]{$a_1(t)$} (out.west);
\draw[->,thick] (a2.south) to[bend left=25] node[right,font=\scriptsize]{$a_2(t)$} (out.east);
\draw[->,thick] (out) -- (env);
\end{tikzpicture}
\caption{Vòng tương tác của hai agent trong một bước thời gian.}
\label{fig:loop}
\end{figure}

\section{Agent A: Q-learning dạng bảng}
Agent RL giữ bảng $Q:\mathcal S\times\mathcal A\to\mathbb R$ gồm $4\times2=8$ giá trị.
Sau mỗi bước, giá trị của cặp trạng thái--hành động vừa dùng được cập nhật
\begin{equation}
Q(s_t,a_t)\leftarrow Q(s_t,a_t)+\alpha\Bigl[r_t+\gamma\max_{a'}Q(s_{t+1},a')-Q(s_t,a_t)\Bigr],
\label{eq:q}
\end{equation}
với $\alpha=0{,}10$, $\gamma=0{,}90$; ở bước cuối episode, mục tiêu chỉ còn $r_t$.
Hành động được chọn theo $\varepsilon$-greedy: ngẫu nhiên với xác suất
$\varepsilon$, ngược lại chọn $\arg\max_aQ(s_t,a)$ (hòa thì chọn ngẫu nhiên).
$\varepsilon$ giảm theo $\varepsilon_e=\max(0{,}02;\;0{,}20\cdot0{,}995^{e})$ ở episode $e$.

\section{Agent B: chính sách dựa trên luật}
\begin{equation}
\pi_{\text{rule}}(d,c)=\begin{cases}0,& d=0,\\ 1,& d=1\text{ và }c=1,\\ 0,& d=1\text{ và }c=0.\end{cases}
\end{equation}
Nghĩa là: nhu cầu thấp thì không giảm giá; nhu cầu cao thì giảm giá chỉ khi đối thủ
vừa giảm giá ở bước trước. Agent này không dùng phần thưởng để thay đổi hành vi.

\section{Phân tích lợi nhuận kỳ vọng một bước}
Trước khi chạy thí nghiệm, có thể tính lợi nhuận kỳ vọng
$\mathbb E[r_i]=N_t\,p_i\,(P(1-\delta a_i)-C)$ cho mọi tổ hợp. Lợi nhuận mỗi sản
phẩm là $10-5=5$ khi không giảm giá và $8-5=3$ khi giảm giá 20\%.

\begin{table}[H]
\centering
\small
\caption{Lợi nhuận kỳ vọng một bước của agent $i$ theo nhu cầu $d$ và hành động hiện tại $b$ của đối thủ (output của chương trình).}
\label{tab:onestep}
\begin{tabular}{ccccc}
\toprule
$d$ & $b=a_{-i}(t)$ & Không giảm giá & Giảm giá & Hành động tốt hơn \\
\midrule
0 & 0 & $10\cdot0{,}30\cdot5=15{,}0$ & $10\cdot0{,}45\cdot3=13{,}5$ & không giảm \\
0 & 1 & $10\cdot0{,}20\cdot5=10{,}0$ & $10\cdot0{,}35\cdot3=10{,}5$ & giảm \\
1 & 0 & $20\cdot0{,}55\cdot5=55{,}0$ & $20\cdot0{,}70\cdot3=42{,}0$ & không giảm \\
1 & 1 & $20\cdot0{,}45\cdot5=45{,}0$ & $20\cdot0{,}60\cdot3=36{,}0$ & không giảm \\
\bottomrule
\end{tabular}
\end{table}

Agent dựa trên luật không bao giờ giảm giá khi nhu cầu thấp, nên dòng thứ hai
(nơi giảm giá có lợi hơn) không xảy ra khi đối thủ là agent luật. Ngoài ra, nếu
agent RL giảm giá ở bước $t$ và nhu cầu ở bước $t+1$ cao, agent luật sẽ giảm giá ở
bước $t+1$, khiến lợi nhuận của RL giảm từ 55 xuống 45. Vì vậy không giảm giá vừa có
lợi trước mắt vừa không kích hoạt phản ứng của đối thủ. Khi đối thủ là agent luật,
chính sách tối ưu là \emph{luôn không giảm giá}, cho lợi nhuận kỳ vọng
$0{,}5\cdot15+0{,}5\cdot55=35{,}0$ mỗi bước. Hàm \texttt{best\_response} trong
chương trình xác nhận điều này bằng cách liệt kê toàn bộ 16 chính sách tất định
trên trạng thái $(d,c)$ (Mục~\ref{sec:exact}).

Một hệ quả quan trọng: hai agent luật đối đầu nhau cũng không bao giờ giảm giá (vì
hành động trước ban đầu là 0), nên agent luật đặt vào ghế của RL đạt đúng 35{,}0.
Trong môi trường tham chiếu, RL \emph{tốt nhất chỉ có thể hòa} agent dựa trên luật.

% ============================================================
\chapter{Cài đặt và thiết lập thí nghiệm}
\label{chap:setup}

\section{Kiến trúc chương trình}
Chương trình phát triển từ mã tham khảo \texttt{agent02\_experiment.py} của đề bài
\cite{complementary}, giữ nguyên mô hình thị trường và các tham số, và thêm các phần
cần cho thí nghiệm A--E và câu 8--10.

\begin{table}[H]
\centering
\small
\caption{Các thành phần chính của chương trình.}
\begin{tabularx}{\textwidth}{>{\ttfamily}p{5.2cm}X}
\toprule
\normalfont\bfseries Thành phần & \bfseries Vai trò \\
\midrule
Config & Tham số mô phỏng, seed, số episode, kích thước cửa sổ \\
ECommerceEnvironment & Sinh nhu cầu, doanh số, phần thưởng và trạng thái tiếp theo \\
QLearningAgent & Q-learning với trạng thái $(d,c)$ như đề bài \\
ExtendedQLearningAgent & Q-learning với trạng thái $(d,c,a_i(t-1))$ (so sánh ở câu 10) \\
ModelBasedAgent & Agent đề xuất ở câu 10: học mô hình rồi lập kế hoạch \\
RuleBasedAgent, TitForTatAgent & Agent luật của đề bài và đối thủ thay thế cho câu 9--10 \\
run\_training, train\_rl\_against\_rule & Vòng huấn luyện chung cho mọi cặp agent \\
evaluate\_pair & Đánh giá tách biệt: đóng băng bảng Q, $\varepsilon=0$, seed riêng \\
exact\_evaluation, best\_response & Giá trị kỳ vọng chính xác của chính sách và chính sách tối ưu \\
run\_baseline, run\_*\_experiment & Thí nghiệm A--E, câu 8, câu 9--10; ghi CSV và hình \\
\bottomrule
\end{tabularx}
\end{table}

\section{Các thay đổi so với mã tham khảo}
Đề bài cho phép sửa mã nếu ghi lại thay đổi và giữ so sánh hợp lệ, tái lập được
\cite{assignment}. Các thay đổi gồm:
\begin{enumerate}[label=\textbf{M\arabic*.}]
\item \textbf{Số ngẫu nhiên chung.} Môi trường dùng một bộ sinh số cho nhu cầu và
  một bộ riêng cho doanh số của từng agent, cùng tách từ một seed bằng
  \texttt{SeedSequence}. Khi đánh giá RL--luật và luật--luật cùng seed, hai lần chạy
  thấy \emph{đúng cùng một chuỗi nhu cầu}; nếu hai agent chọn cùng hành động thì
  doanh số cũng giống hệt. Nhờ vậy chênh lệch đo được đến từ chính sách, không
  phải từ may rủi của chuỗi nhu cầu.
\item \textbf{Giao diện agent thống nhất} (\texttt{observe}, \texttt{act},
  \texttt{learn}). Cùng một vòng huấn luyện và đánh giá dùng được cho RL--luật,
  RL--RL và các agent mới. Mỗi agent tự tạo quan sát từ trạng thái chung
  $(d_t,a_1(t-1),a_2(t-1))$.
\item \textbf{Đánh giá chính xác.} Hàm \texttt{exact\_evaluation} lan truyền phân
  phối của trạng thái chung qua 20 bước và tính lợi nhuận kỳ vọng của chính sách
  tham lam mà không cần lấy mẫu. Hàm \texttt{best\_response} liệt kê mọi chính sách
  tất định để tìm mức tối ưu cho từng cách biểu diễn trạng thái.
\item \textbf{Thí nghiệm bổ sung:} đường học theo cửa sổ, $\varepsilon$ cố định,
  10 seed độc lập, hai agent cùng học, thay đổi mức giảm giá, hàm thưởng và đối
  thủ, agent cải tiến.
\end{enumerate}

\section{Tham số và quy trình đánh giá}
\begin{table}[H]
\centering
\small
\caption{Tham số tham chiếu và cấu hình tái lập.}
\label{tab:params}
\begin{tabular}{ll}
\toprule
Tham số & Giá trị \\
\midrule
Số agent & 2 \\
Số episode huấn luyện $\times$ số bước & $5.000\times20=100.000$ bước \\
Tốc độ học $\alpha$, hệ số chiết khấu $\gamma$ & 0{,}10; 0{,}90 \\
Khám phá $\varepsilon$ & $0{,}20\to0{,}02$, hệ số giảm 0{,}995 mỗi episode \\
Giá gốc, giá vốn, mức giảm & 10; 5; 20\% \\
Seed baseline (môi trường) & 42; seed của agent RL $=42+10.000$ \\
Seed đánh giá baseline & $42+500.000$ \\
Seed nhiều lần chạy & 11, 22, 33, 44, 55, 66, 77, 88, 99, 110 \\
Seed đánh giá nhiều lần chạy & seed huấn luyện $+800.000$ \\
Đánh giá & 500 episode $\times$ 20 bước $=10.000$ bước, $\varepsilon=0$, bảng Q đóng băng \\
Môi trường phần mềm & Python 3.12.13, NumPy 2.5.2, Matplotlib 3.11.1 \\
Thời gian chạy toàn bộ & khoảng 40 giây trên MacBook (CPU) \\
\bottomrule
\end{tabular}
\end{table}

Các chỉ số đo: tổng thưởng, thưởng trung bình mỗi bước
$\bar R_i=\frac1T\sum_{t=1}^{T}r_i(t)$, số lần giảm/không giảm giá, tỷ lệ giảm giá
$\rho_i$, bảng Q và chính sách tham lam. Với nhiều seed, báo cáo dùng trung bình và
độ lệch chuẩn mẫu ($n-1$). Chạy lại chương trình hai lần cho nhật ký giống hệt
nhau, nghĩa là kết quả tái lập hoàn toàn với cùng môi trường phần mềm.

Chuỗi chính sách như \texttt{0100} liệt kê hành động tham lam ở các trạng thái
theo thứ tự $(0,0),(0,1),(1,0),(1,1)$ của $(d,c)$; 1 là giảm giá.

\section{Trích đoạn mã nguồn chính}
\subsection{Mô hình thị trường và môi trường}
Hàm xác suất chuyển đổi cài đặt công thức~\eqref{eq:conv}; \texttt{step} nhận hai
hành động đồng thời, trả phần thưởng~\eqref{eq:reward} rồi sinh nhu cầu mới độc lập.
\CodeLines{107}{135}
\CodeLines{155}{211}

\subsection{Agent Q-learning}
\texttt{act} cài đặt $\varepsilon$-greedy với phá hòa ngẫu nhiên; \texttt{learn}
cài đặt công thức~\eqref{eq:q}.
\CodeLines{264}{310}

\subsection{Agent dựa trên luật}
\CodeLines{409}{446}

\subsection{Vòng huấn luyện}
Hai agent quan sát và chọn hành động trước khi môi trường thay đổi; chỉ agent có
\texttt{learns = True} cập nhật bảng Q.
\CodeLines{590}{628}

\subsection{Đánh giá tách biệt và đánh giá chính xác}
\label{sec:exact}
\texttt{evaluate\_pair} chạy chính sách đã đóng băng với $\varepsilon=0$ trên seed
riêng. \texttt{exact\_evaluation} tính kỳ vọng chính xác: ở mỗi bước, với mỗi trạng
thái chung và mỗi tổ hợp hành động, cộng thưởng kỳ vọng $N_t p_i\cdot m_i$ ($m_i$ là biên
lợi nhuận) có trọng số xác suất, rồi chuyển khối xác suất sang trạng thái kế tiếp.
\CodeLines{686}{756}
\CodeLines{477}{566}

% ============================================================
\chapter{Kết quả thí nghiệm}
\label{chap:results}

\section{Thí nghiệm A: baseline}
Agent RL được huấn luyện 5.000 episode với seed 42 trước agent luật, sau đó được
đánh giá riêng trong 500 episode. Để so sánh công bằng, agent luật cũng được đặt vào
đúng ghế của RL và chạy với cùng seed đánh giá.

\begin{Output}
Separate evaluation results (500 episodes, seed 500042):
RL vs Rule    | RL (Q-learning)  | total =  345565.0 | mean reward/step = 34.5565 | discount =     0 | no discount = 10000
RL vs Rule    | Rule-based       | total =  346720.0 | mean reward/step = 34.6720 | discount =     0 | no discount = 10000
Rule vs Rule  | Rule-based       | total =  345565.0 | mean reward/step = 34.5565 | discount =     0 | no discount = 10000
Rule vs Rule  | Rule-based       | total =  346720.0 | mean reward/step = 34.6720 | discount =     0 | no discount = 10000

Expected reward/step (exact): RL greedy = 35.0000, rule in seat 1 = 35.0000,
best (d, c) policy = 35.0000 {(0, 0): 0, (0, 1): 0, (1, 0): 0, (1, 1): 0}
\end{Output}

\begin{table}[H]
\centering
\small
\caption{Kết quả đánh giá baseline (seed 42, 10.000 bước đánh giá).}
\label{tab:baseline}
\begin{tabular}{llrrrr}
\toprule
Lần chạy & Agent & Tổng thưởng & TB/bước & Giảm giá & Không giảm \\
\midrule
RL vs Luật & RL (ghế 1) & 345.565 & 34{,}5565 & 0 & 10.000 \\
RL vs Luật & Luật (ghế 2) & 346.720 & 34{,}6720 & 0 & 10.000 \\
Luật vs Luật & Luật (ghế 1) & 345.565 & 34{,}5565 & 0 & 10.000 \\
Luật vs Luật & Luật (ghế 2) & 346.720 & 34{,}6720 & 0 & 10.000 \\
\bottomrule
\end{tabular}
\end{table}

Chính sách tham lam học được là \emph{không giảm giá ở cả bốn trạng thái}. Agent
luật cũng không giảm giá lần nào, nên hai lần chạy cho kết quả trùng đến từng đơn vị
nhờ dùng chung số ngẫu nhiên (M1). Chênh lệch giữa ghế 1 (34{,}56) và ghế 2 (34{,}67)
chỉ do doanh số ngẫu nhiên khác nhau; giá trị kỳ vọng của cả hai là 35{,}00.

\begin{Output}
Q-table: Q[d, c, action]; Action 0 = no discount; Action 1 = discount
State=(0, 0) | Q0=209.034622 | Q1=200.012813 | Greedy action=0
State=(0, 1) | Q0=227.101684 | Q1=181.167889 | Greedy action=0
State=(1, 0) | Q0=266.318271 | Q1=228.504950 | Greedy action=0
State=(1, 1) | Q0=228.334394 | Q1=219.318896 | Greedy action=0
\end{Output}

Mọi trạng thái có $Q_0>Q_1$ nên chính sách tham lam chọn không giảm giá. Độ chênh
nhỏ nhất (khoảng 9) nằm ở $(0,0)$ và $(1,1)$. Ở $(0,0)$, điều này khớp với
Bảng~\ref{tab:onestep}: khi nhu cầu thấp, giảm giá chỉ thiệt 1{,}5 mỗi bước. Trạng
thái $(1,1)$ ít được thăm (đối thủ chỉ giảm giá sau khi RL giảm giá) nên ước lượng
của nó dựa trên ít dữ liệu. Giá trị Q vào khoảng 200--270 vì nó
ước lượng tổng thưởng chiết khấu của nhiều bước còn lại, gần
$35/(1-0{,}9)=350$ ở đầu episode và chỉ còn khoảng 35 ở bước cuối.

\section{Thí nghiệm B: quá trình học}
Bảng~\ref{tab:windows} tính $\bar R_{i,k}=\sum_{t\in W_k}r_i(t)/(20K)$ với cửa sổ
$K=500$ episode; Hình~\ref{fig:train} dùng $K=100$.

\begin{table}[H]
\centering
\small
\caption{Thưởng trung bình mỗi bước trong huấn luyện theo cửa sổ 500 episode (seed 42).}
\label{tab:windows}
\begin{tabular}{lrrrr}
\toprule
Episode & RL & Luật & Tỷ lệ giảm giá của RL & $\varepsilon$ đầu cửa sổ \\
\midrule
1--500 & 33{,}474 & 32{,}862 & 0{,}193 & 0{,}200 \\
501--1000 & 33{,}412 & 33{,}172 & 0{,}139 & 0{,}020 \\
1001--1500 & 33{,}382 & 32{,}620 & 0{,}195 & 0{,}020 \\
1501--2000 & 33{,}541 & 32{,}922 & 0{,}201 & 0{,}020 \\
2001--2500 & 33{,}015 & 32{,}439 & 0{,}201 & 0{,}020 \\
2501--3000 & 33{,}965 & 33{,}284 & 0{,}170 & 0{,}020 \\
3001--3500 & 33{,}841 & 33{,}472 & 0{,}146 & 0{,}020 \\
3501--4000 & 32{,}492 & 31{,}774 & 0{,}232 & 0{,}020 \\
4001--4500 & 33{,}321 & 32{,}849 & 0{,}179 & 0{,}020 \\
4501--5000 & 33{,}713 & 33{,}177 & 0{,}131 & 0{,}020 \\
\bottomrule
\end{tabular}
\end{table}

\Fig{episode_rewards.png}{Thí nghiệm B: thưởng mỗi bước theo episode, trung bình cửa sổ $K=100$, tỷ lệ giảm giá của RL và $\varepsilon$ (seed 42).\label{fig:train}}

\textbf{Quan sát.} Thưởng huấn luyện của RL dao động quanh 33--34 suốt 5.000 episode
và không tăng rõ rệt sau vài trăm episode đầu. Tỷ lệ giảm giá trong huấn luyện vẫn ở
mức 0{,}13--0{,}23 dù $\varepsilon$ đã chạm sàn 0{,}02 từ khoảng episode 460. Phần
khám phá ngẫu nhiên chỉ giải thích được khoảng 0{,}01 trong tỷ lệ này, nên phần lớn
các lần giảm giá đến từ chính sách tham lam tạm thời chọn giảm giá.
Hình~\ref{fig:pv} xác nhận điều đó: giá trị kỳ vọng chính xác của chính sách tham lam,
đo mỗi 25 episode, liên tục nhảy giữa 35{,}0 (tối ưu) và các mức 27--33.

\Fig{policy_value_during_training.png}{Giá trị kỳ vọng chính xác của chính sách tham lam trong lúc huấn luyện (seed 42).\label{fig:pv}}

\Fig{q_value_trajectories.png}{Quỹ đạo tám giá trị Q cuối mỗi episode (seed 42).\label{fig:qtraj}}

\textbf{Giải thích.} Hình~\ref{fig:qtraj} cho thấy các giá trị Q không hội tụ mà dao
động với biên độ vài chục đơn vị, lớn hơn khoảng cách giữa $Q_0$ và $Q_1$ ở trạng
thái $(0,0)$. Có ba nguyên nhân hợp lý: (i) doanh số là biến nhị thức nên thưởng
mỗi bước có độ lệch chuẩn lớn (khoảng 11 khi nhu cầu cao); (ii) $\alpha=0{,}1$
cố định nên mỗi cập nhật chỉ làm trung bình trên khoảng 10 mẫu gần nhất; (iii) trạng
thái không chứa vị trí bước trong episode, nên cùng một trạng thái có mục tiêu khoảng
$300$ ở đầu episode nhưng chỉ $r_t\approx35$ ở bước cuối. Vì vậy, thưởng huấn luyện
tăng hay giảm \emph{không} đủ để kết luận về hội tụ: trong lần chạy này, thưởng
huấn luyện gần như không đổi, nhưng chính sách cuối cùng vẫn tối ưu. Chính sách
cuối cùng tối ưu một phần nhờ may mắn: nếu dừng ở episode khác, chính sách có thể
kém hơn.

\section{Thí nghiệm C: tốc độ học}
Ba cấu hình $\alpha\in\{0{,}01;0{,}10;0{,}50\}$ được chạy với 10 seed; các tham số
khác giữ nguyên.

\begin{table}[H]
\centering
\small
\caption{Ảnh hưởng của tốc độ học (10 seed). ``Tối ưu'' là số seed có chính sách cuối đạt 35{,}00 kỳ vọng; ``Đổi chính sách'' là số lần chính sách tham lam thay đổi trong 2.500 episode cuối.}
\label{tab:lr}
\footnotesize
\begin{tabular}{lrrrrrr}
\toprule
$\alpha$ & Đánh giá (TB $\pm$ SD) & Giảm giá & Huấn luyện & 1.000 ep. cuối & Đổi chính sách & Tối ưu \\
\midrule
0{,}01 & $34{,}546\pm1{,}328$ & 0{,}049 & 32{,}245 & 34{,}414 & 17{,}1 & 9/10 \\
0{,}10 & $34{,}724\pm0{,}747$ & 0{,}040 & 33{,}600 & 33{,}566 & 256{,}3 & 9/10 \\
0{,}50 & $32{,}912\pm2{,}457$ & 0{,}259 & 31{,}876 & 31{,}763 & 827{,}7 & 5/10 \\
\bottomrule
\end{tabular}
\end{table}

\Fig{learning_rate_curves.png}{Thí nghiệm C: đường học (TB $\pm$ SD trên 10 seed, cửa sổ 100 episode).\label{fig:lrcurve}}
\Fig{learning_rate_q_trajectory.png}{Thí nghiệm C: hiệu $Q((0,0),0)-Q((0,0),1)$ theo episode, một lần chạy với seed 1042 cho mỗi $\alpha$.\label{fig:lrq}}
\Fig[0.8]{learning_rate_comparison.png}{Thí nghiệm C: thưởng đánh giá theo $\alpha$.}

\section{Thí nghiệm D: mức khám phá}
Năm cấu hình được so sánh: lịch tham chiếu $0{,}20\to0{,}02$; $\varepsilon$ cố định
0; 0{,}05; 0{,}30; và lịch cao $0{,}50\to0{,}05$ của tài liệu code
\cite{complementary}. Cấu hình $\varepsilon=0$ được thêm để xem điều gì xảy ra khi
agent ngừng khám phá ngay từ đầu.

\begin{table}[H]
\centering
\small
\caption{Ảnh hưởng của khám phá (10 seed). Thưởng huấn luyện là trung bình 5.000 episode.}
\label{tab:explore}
\footnotesize
\begin{tabular}{lrrrrrr}
\toprule
Cấu hình & Đánh giá (TB $\pm$ SD) & Giảm giá (ĐG) & Huấn luyện & Giảm giá (HL) & Đổi CS & Tối ưu \\
\midrule
Lịch $0{,}20\to0{,}02$ & $34{,}724\pm0{,}747$ & 0{,}040 & 33{,}600 & 0{,}167 & 256{,}3 & 9/10 \\
Cố định 0{,}00 & $29{,}719\pm3{,}852$ & 0{,}415 & 29{,}720 & 0{,}414 & 0{,}0 & 3/10 \\
Cố định 0{,}05 & $33{,}775\pm2{,}267$ & 0{,}130 & 32{,}907 & 0{,}231 & 502{,}3 & 7/10 \\
Cố định 0{,}30 & $29{,}476\pm3{,}104$ & 0{,}461 & 30{,}910 & 0{,}400 & 1.217{,}0 & 2/10 \\
Lịch $0{,}50\to0{,}05$ & $32{,}411\pm2{,}656$ & 0{,}300 & 32{,}867 & 0{,}236 & 505{,}2 & 4/10 \\
\bottomrule
\end{tabular}
\end{table}

\Fig{exploration_curves.png}{Thí nghiệm D: đường học theo cấu hình khám phá (TB $\pm$ SD trên 10 seed).\label{fig:expcurve}}
\Fig[0.8]{exploration_comparison.png}{Thí nghiệm D: thưởng đánh giá theo cấu hình khám phá.}

\section{Thí nghiệm E: so sánh công bằng trên nhiều seed}
RL (lịch tham chiếu) và agent luật được đánh giá trong các lần chạy tách biệt, cùng
ghế 1, cùng đối thủ là agent luật và cùng seed đánh giá (seed huấn luyện $+800.000$).

\begin{table}[H]
\centering
\small
\caption{Bảng kết quả theo gợi ý của đề bài (đánh giá tham lam, 10 seed, 10.000 bước mỗi seed).}
\label{tab:suggested}
\begin{tabular}{lrrr}
\toprule
Agent / cấu hình & Thưởng TB & Độ lệch chuẩn & Tỷ lệ giảm giá \\
\midrule
RL, baseline ($\alpha=0{,}10$, lịch $0{,}20\to0{,}02$) & 34{,}724 & 0{,}747 & 0{,}040 \\
RL, tốc độ học thấp ($\alpha=0{,}01$) & 34{,}546 & 1{,}328 & 0{,}049 \\
RL, tốc độ học cao ($\alpha=0{,}50$) & 32{,}912 & 2{,}457 & 0{,}259 \\
RL, khám phá khác ($\varepsilon=0{,}05$ cố định) & 33{,}775 & 2{,}267 & 0{,}130 \\
RL, khám phá khác ($\varepsilon=0{,}30$ cố định) & 29{,}476 & 3{,}104 & 0{,}461 \\
Agent dựa trên luật & 34{,}985 & 0{,}178 & 0{,}000 \\
\bottomrule
\end{tabular}
\end{table}

\begin{table}[H]
\centering
\small
\caption{So sánh từng cặp RL--luật trên cùng chuỗi nhu cầu (\texttt{multi\_seed\_summary.csv}).}
\label{tab:paired}
\begin{tabular}{rrrrc}
\toprule
Seed & RL & Luật & RL $-$ Luật & Chính sách RL \\
\midrule
11 & 35{,}189 & 35{,}189 & 0{,}000 & 0100 \\
22 & 34{,}881 & 34{,}881 & 0{,}000 & 0100 \\
33 & 32{,}641 & 35{,}258 & $-2{,}617$ & 1000 \\
44 & 34{,}931 & 34{,}931 & 0{,}000 & 0000 \\
55 & 34{,}951 & 34{,}951 & 0{,}000 & 0000 \\
66 & 34{,}765 & 34{,}765 & 0{,}000 & 0100 \\
77 & 34{,}904 & 34{,}904 & 0{,}000 & 0100 \\
88 & 34{,}780 & 34{,}780 & 0{,}000 & 0000 \\
99 & 34{,}972 & 34{,}972 & 0{,}000 & 0000 \\
110 & 35{,}226 & 35{,}226 & 0{,}000 & 0100 \\
\bottomrule
\end{tabular}
\end{table}

Ở 9/10 seed, RL và luật cho kết quả \emph{giống hệt}. Các chính sách \texttt{0100}
chọn giảm giá ở trạng thái $(0,1)$, nhưng trạng thái này không bao giờ xuất hiện khi
RL không giảm giá: agent luật chỉ giảm giá sau khi RL giảm giá. Vì vậy hành vi thực
tế của chúng vẫn là không bao giờ giảm giá. Seed 33 học sai ở trạng thái $(0,0)$:
giảm giá khi nhu cầu thấp, làm lợi nhuận giảm còn 32{,}64. Hiệu cặp RL$-$luật có
trung bình $-0{,}262$, độ lệch chuẩn $0{,}828$, khoảng tin cậy 95\% (phân phối $t$,
9 bậc tự do) là $[-0{,}854;\;0{,}330]$, có chứa 0.

\Fig{multi_seed_comparison.png}{Thí nghiệm E: RL và luật trên cùng chuỗi nhu cầu theo từng seed.}

% ============================================================
\chapter{Trả lời mười câu hỏi}
\label{chap:questions}

\needspace{10\baselineskip}
\section*{Câu 1. Mô hình hóa agent và môi trường}
\addcontentsline{toc}{section}{Câu 1. Mô hình hóa agent và môi trường}
\begin{itemize}
\item \textbf{Agent:} hai người bán. Agent 1 là agent Q-learning, agent 2 là agent
  dựa trên luật (ở câu 8, cả hai đều là agent Q-learning).
\item \textbf{Môi trường:} thị trường thương mại điện tử gồm nhu cầu $d_t$, số khách
  $N_t$, xác suất mua~\eqref{eq:conv} và cơ chế bán theo phân phối nhị thức.
\item \textbf{Không gian trạng thái} của mỗi agent: $s_i=(d,c)\in\{0,1\}^2$, gồm 4
  trạng thái. Trạng thái đầy đủ của môi trường là $(d_t,a_1(t-1),a_2(t-1))$, gồm 8
  trạng thái; mỗi agent chỉ thấy một phần.
\item \textbf{Không gian hành động:} $\mathcal A=\{0,1\}$ (không giảm/giảm giá),
  chọn đồng thời.
\item \textbf{Hàm thưởng:} lợi nhuận~\eqref{eq:reward}, ngẫu nhiên vì $S_i$ là
  biến nhị thức.
\item \textbf{Chuyển trạng thái:} $d_{t+1}$ đều trên $\{0,1\}$, độc lập với hành
  động; $c_i(t+1)=a_{-i}(t)$. Mỗi episode có 20 bước và đặt lại hành động trước về 0.
\end{itemize}
Đây là hệ \emph{nhiều agent} vì phần thưởng và trạng thái của mỗi agent phụ thuộc vào
hành động của agent kia: $p_i$ giảm 0{,}10 khi đối thủ giảm giá, và $c_i(t+1)$ chính
là hành động của đối thủ. Nếu coi đối thủ là một phần của môi trường thì môi trường
không còn cố định với agent RL khi đối thủ cũng học (câu 8), và ngay cả với đối thủ
luật, môi trường phản ứng với hành động của RL (giảm giá kéo theo đối thủ giảm giá).
Mô hình một agent sẽ bỏ qua cả tương tác chiến lược này lẫn câu hỏi về lợi ích của
từng bên.

\needspace{10\baselineskip}
\section*{Câu 2. Quyết định dựa trên học và dựa trên luật}
\addcontentsline{toc}{section}{Câu 2. Quyết định dựa trên học và dựa trên luật}
\begin{table}[H]
\centering
\small
\caption{So sánh agent RL và agent dựa trên luật.}
\begin{tabularx}{\textwidth}{>{\bfseries}p{2.8cm}XX}
\toprule
Tiêu chí & Agent RL (Q-learning) & Agent dựa trên luật \\
\midrule
Ra quyết định & $\arg\max_aQ(s,a)$, có khám phá khi huấn luyện & Tra bảng if--else cố định \\
Dùng kinh nghiệm & Cập nhật $Q$ sau mỗi bước từ thưởng quan sát được & Không dùng thưởng; tri thức nằm trong luật do người viết \\
Khả năng thích ứng & Tự đổi chính sách khi thị trường đổi; ở mức giảm 5\% đạt 36{,}72 so với 34{,}99 của luật & Phải viết lại luật; ở mức giảm 5\% vẫn giữ nguyên hành vi \\
Độ ổn định & Phụ thuộc seed, $\alpha$, $\varepsilon$; SD 0{,}75 (baseline) & Hoàn toàn nhất quán; SD 0{,}18 chỉ do nhiễu doanh số \\
Tính toán & 100.000 bước huấn luyện ($\approx0{,}5$ s), bảng 8 giá trị & Không cần huấn luyện, $O(1)$ mỗi quyết định \\
Giải thích được & Cần đọc bảng Q; khó hơn khi trạng thái lớn & Đọc trực tiếp từ luật \\
\bottomrule
\end{tabularx}
\end{table}
\textbf{Khi nên dùng agent luật:} khi có chính sách được kiểm chứng hoặc ràng buộc
pháp lý rõ ràng, ví dụ quy định không được bán dưới giá vốn hay chương trình khuyến
mãi theo lịch cố định trong dịp lễ. Khi đó tính nhất quán và khả năng giải thích
quan trọng hơn vài phần trăm lợi nhuận. Trong thí nghiệm này, luật tham chiếu tình cờ
là chính sách tối ưu nên RL không thể làm tốt hơn.
\textbf{Khi nên dùng agent RL:} khi điều kiện thị trường thay đổi và không ai biết
trước luật tốt nhất, ví dụ chi phí, độ nhạy giá hoặc mức giảm giá thay đổi theo mùa.
Kịch bản giảm 5\% ở câu 9 cho thấy RL tự phát hiện rằng luôn giảm giá có lợi hơn,
trong khi luật cũ mất khoảng 2{,}35 lợi nhuận mỗi bước.

\needspace{10\baselineskip}
\section*{Câu 3. Cập nhật Q-learning}
\addcontentsline{toc}{section}{Câu 3. Cập nhật Q-learning}
Công thức~\eqref{eq:q} dịch $Q(s,a)$ một phần $\alpha$ về phía \emph{mục tiêu TD}
$r+\gamma\max_{a'}Q(s',a')$. Mục tiêu này là ước lượng một bước của tổng thưởng
chiết khấu: thưởng nhận ngay cộng giá trị tốt nhất có thể đạt từ trạng thái kế tiếp.
Hiệu $\delta=\text{mục tiêu}-Q(s,a)$ là \emph{sai số TD}: dương thì giá trị cũ đang
bị đánh giá thấp. $\alpha$ quyết định mức tin vào mẫu mới; $\gamma$ quyết định trọng
số của tương lai. Toán tử $\max$ làm Q-learning học \emph{off-policy}: nó ước lượng
giá trị của chính sách tham lam dù hành động đang thực hiện có thể là khám phá.

Với $Q(s,a)=4$, $\alpha=0{,}1$, $r=6$, $\gamma=0{,}9$, $\max_{a'}Q(s',a')=5$:
\begin{align*}
\text{Mục tiêu TD} &= r+\gamma\max_{a'}Q(s',a')=6+0{,}9\times5=6+4{,}5=10{,}5,\\
\text{Sai số TD } \delta &= 10{,}5-4=6{,}5,\\
\alpha\,\delta &= 0{,}1\times6{,}5=0{,}65,\\
Q_{\text{mới}}(s,a) &= 4+0{,}65=4{,}65.
\end{align*}
Chương trình in cùng kết quả: \texttt{Q3: target = 6.0 + 0.9 * 5.0 = 10.50; TD
error = 6.50; new Q = 4.0 + 0.1 * 6.50 = 4.65}. \textbf{Diễn giải:} kinh nghiệm mới
cho thấy hành động $a$ tốt hơn mức agent đang tin (10{,}5 so với 4), nhưng với
$\alpha=0{,}1$ agent chỉ dịch 10\% quãng đường, từ 4 lên 4{,}65. Nếu cùng mục tiêu
lặp lại, $Q$ tiến dần về 10{,}5 theo cấp số nhân với hệ số $(1-\alpha)$; một mẫu
nhiễu đơn lẻ không làm thay đổi đánh giá quá mạnh.

\needspace{10\baselineskip}
\section*{Câu 4. Khám phá và khai thác}
\addcontentsline{toc}{section}{Câu 4. Khám phá và khai thác}
$\varepsilon$-greedy chọn hành động ngẫu nhiên với xác suất $\varepsilon$ để tiếp
tục thu thập thông tin về những hành động có giá trị chưa chắc chắn; phần còn lại
khai thác hành động có $Q$ cao nhất. Bảng~\ref{tab:explore} và
Hình~\ref{fig:expcurve} cho thấy:
\begin{itemize}
\item \textbf{Ngừng khám phá quá sớm} ($\varepsilon=0$): đánh giá $29{,}72\pm3{,}85$,
  chỉ 3/10 seed tối ưu, chính sách không đổi lần nào trong nửa sau. Lý do: mọi
  lợi nhuận đều dương và $Q$ khởi tạo bằng 0, nên hành động được thử đầu tiên ở mỗi
  trạng thái (chọn ngẫu nhiên do hòa) có ngay $Q>0$ và được chọn mãi mãi. Agent bị
  khóa vào lựa chọn ngẫu nhiên ban đầu; tỷ lệ giảm giá 0{,}415 gần với kết quả tung
  đồng xu.
\item \textbf{Khám phá quá nhiều} ($\varepsilon=0{,}30$ cố định): thưởng huấn luyện
  chỉ 30{,}91 vì khoảng 15\% bước là giảm giá ngẫu nhiên, mỗi lần làm mất lợi nhuận
  và kích hoạt phản ứng giảm giá của đối thủ. Đáng chú ý, chính sách tham lam cuối
  cùng cũng kém (29{,}48, 2/10 tối ưu). Một giải thích phù hợp với dữ liệu: trạng
  thái $(d,c)$ không chứa hành động trước của chính RL, nên giá trị Q mà Q-learning
  hội tụ tới phụ thuộc vào tần suất RL tự giảm giá trong lúc huấn luyện. Khám phá
  nhiều làm thay đổi phân phối đó. Thí nghiệm hiện tại chưa đủ để chứng minh cơ chế
  này; cần một thí nghiệm riêng với trạng thái đầy đủ.
\item \textbf{Lịch giảm dần $0{,}20\to0{,}02$} cho kết quả tốt nhất ($34{,}72$,
  9/10), vì khám phá nhiều lúc bảng Q còn trống, rồi khai thác khi đã có ước lượng.
  $\varepsilon=0{,}05$ cố định đứng thứ hai ($33{,}78$, 7/10); lịch cao
  $0{,}50\to0{,}05$ kém hơn ($32{,}41$, 4/10).
\end{itemize}
Thưởng huấn luyện và thưởng đánh giá đo hai thứ khác nhau: với $\varepsilon=0{,}30$,
thưởng huấn luyện (30{,}91) cao hơn thưởng đánh giá (29{,}48) dù đánh giá không có
hành động ngẫu nhiên. Vì vậy chỉ thưởng đánh giá mới phản ánh chất lượng chính sách.

\needspace{10\baselineskip}
\section*{Câu 5. Phần thưởng và mục tiêu kinh doanh}
\addcontentsline{toc}{section}{Câu 5. Phần thưởng và mục tiêu kinh doanh}
Giảm giá làm tăng xác suất mua thêm 0{,}15 (công thức~\eqref{eq:conv}), nhưng giá
bán giảm từ 10 xuống $10\times(1-0{,}20)=8$ trong khi giá vốn vẫn là 5:
\begin{align*}
\text{Không giảm giá: } & 10-5=5 \text{ trên mỗi sản phẩm},\\
\text{Giảm giá 20\%: } & 10\times0{,}8-5=8-5=3 \text{ trên mỗi sản phẩm}.
\end{align*}
Biên lợi nhuận giảm 40\%, nên giảm giá chỉ có lợi khi số đơn tăng hơn
$5/3\approx1{,}67$ lần. Khi nhu cầu thấp, tỷ lệ tăng là $0{,}45/0{,}30=1{,}5$; khi
nhu cầu cao là $0{,}70/0{,}55\approx1{,}27$. Cả hai đều dưới 1{,}67, nên giảm giá
làm giảm lợi nhuận kỳ vọng (Bảng~\ref{tab:onestep}).

\textbf{Thay đổi hàm thưởng sẽ thay đổi hành vi học được.} Chương trình chạy thêm
kịch bản phần thưởng là \emph{doanh thu} $S_i\cdot P(1-\delta a_i)$ thay cho lợi
nhuận (Bảng~\ref{tab:scenario}). Với doanh thu, giảm giá khi nhu cầu thấp cho
$10\times0{,}45\times8=36$ so với $10\times0{,}30\times10=30$, tức có lợi trước mắt.
Agent RL tham chiếu nâng tỷ lệ giảm giá từ 0{,}04 lên 0{,}50, và lợi nhuận thực tế
giảm từ 34{,}72 xuống $29{,}07\pm3{,}20$. Agent tối ưu doanh thu nhưng làm doanh
nghiệp mất lợi nhuận. Một quan sát phụ: lợi nhuận mỗi đơn với mức giảm 10\% là
$(5;4)$, đúng bằng một nửa doanh thu mỗi đơn với mức giảm 20\% là $(10;8)$. Vì
Q-learning không đổi khi nhân thưởng với hằng số dương, hai kịch bản này cho cùng
chính sách ở mọi seed (cùng tỷ lệ giảm giá 0{,}500 và cùng 338{,}5 lần đổi chính
sách). Các thay đổi khác cũng làm đổi hành vi: phạt mỗi lần giảm giá sẽ đẩy agent
về phía không giảm giá; thêm thưởng cho thị phần hay khách hàng trung thành có thể
làm giảm giá trở nên có lợi. Hàm thưởng phải phản ánh đúng mục tiêu kinh doanh, vì
agent sẽ tối ưu chính xác thứ được đo.

\needspace{10\baselineskip}
\section*{Câu 6. Đánh giá hiệu năng}
\addcontentsline{toc}{section}{Câu 6. Đánh giá hiệu năng}
Thưởng trung bình mỗi bước $\bar R_i=\frac1T\sum_t r_i(t)$ trong đánh giá tách biệt:
\begin{itemize}
\item Baseline seed 42, $T=10.000$: RL $345.565/10.000=34{,}557$; agent luật ở cùng
  ghế cũng $34{,}557$; agent luật ở ghế đối thủ $34{,}672$ (Bảng~\ref{tab:baseline}).
\item 10 seed: RL $34{,}724\pm0{,}747$; luật $34{,}985\pm0{,}178$
  (Bảng~\ref{tab:suggested}); giá trị kỳ vọng chính xác lần lượt
  $34{,}745\pm0{,}807$ và $35{,}000$.
\end{itemize}
Một lần chạy không đủ để kết luận một phương pháp tốt hơn vì:
(i) \emph{nhiễu môi trường}: ngay cả khi chính sách giống hệt, ghế 1 và ghế 2 lệch
nhau 0{,}12 chỉ do doanh số ngẫu nhiên, và giữa các seed đánh giá, cùng một chính
sách ``không giảm giá'' cho kết quả từ 34{,}77 đến 35{,}26;
(ii) \emph{nhiễu huấn luyện}: chính sách cuối phụ thuộc seed; 9 seed hòa với luật
nhưng seed 33 thua 2{,}62. Nếu chỉ chạy seed 33, ta sẽ kết luận RL kém hẳn; nếu chỉ
chạy seed 11 thì kết luận ngược lại; (iii) \emph{phụ thuộc cấu hình}: kết quả đổi theo
$\alpha$, $\varepsilon$ và tham số thị trường (câu 9). Với 10 seed và so sánh từng
cặp, khoảng tin cậy 95\% của hiệu RL$-$luật là $[-0{,}854;0{,}330]$, tức không có
bằng chứng một bên tốt hơn. Đây cũng chỉ là kết luận trong mô hình mô phỏng này,
không phải cho thị trường thật.

\needspace{10\baselineskip}
\section*{Câu 7. Tốc độ học và độ ổn định}
\addcontentsline{toc}{section}{Câu 7. Tốc độ học và độ ổn định}
Theo Bảng~\ref{tab:lr} và các Hình~\ref{fig:lrcurve}, \ref{fig:lrq}:
\begin{itemize}
\item \textbf{$\alpha=0{,}01$ (chậm, ổn định):} mỗi mẫu chỉ đóng góp 1\%, nên ước
  lượng ít nhiễu (chỉ 17{,}1 lần đổi chính sách trong nửa sau, so với 256{,}3 ở
  $\alpha=0{,}1$). Đổi lại, agent học rất chậm: đường học trung bình đứng quanh 30{,}6
  trong khoảng 2.500 episode đầu, sau đó mới tăng lên khoảng 34{,}8. Vì $Q$ khởi tạo
  bằng 0 và mọi thưởng đều dương, hành động được thử nhiều hơn lúc đầu sẽ có $Q$ cao
  hơn trong thời gian dài. Hình~\ref{fig:lrq} cho thấy ở seed 1042, hiệu
  $Q_0-Q_1$ tại $(0,0)$ vẫn âm suốt 5.000 episode (dù đang tăng dần về 0), tức agent
  chưa kịp sửa lựa chọn ban đầu. Đánh giá $34{,}55\pm1{,}33$; 9/10 seed tối ưu.
\item \textbf{$\alpha=0{,}10$ (tham chiếu):} học nhanh trong vài trăm episode, kết quả
  tốt nhất ($34{,}72\pm0{,}75$; 9/10), nhưng hiệu $Q_0-Q_1$ vẫn đổi dấu 23\% thời
  gian nửa sau ở seed 1042.
\item \textbf{$\alpha=0{,}50$ (nhanh, dao động):} mỗi cập nhật thay một nửa giá trị
  cũ bằng một mẫu nhiễu, nên $Q$ gần như chỉ phản ánh vài mẫu gần nhất. Hiệu
  $Q_0-Q_1$ dao động trong khoảng $\pm200$ và âm 50\% thời gian; chính sách đổi 827{,}7
  lần trong nửa sau. Thưởng đánh giá thấp nhất ($32{,}91\pm2{,}46$, chỉ 5/10 tối ưu).
\end{itemize}
Tóm lại, $\alpha$ lớn làm agent nhạy với quan sát mới nhưng ước lượng nhiễu và chính
sách không ổn định; $\alpha$ nhỏ làm ước lượng mượt nhưng cần nhiều dữ liệu hơn.
Với thưởng nhiều nhiễu như ở đây, $\alpha$ giảm dần theo số lần thăm (ví dụ
$1/N(s,a)$) hoặc ước lượng mô hình bằng trung bình mẫu (câu 10) là lựa chọn tốt hơn
$\alpha$ cố định.

\needspace{10\baselineskip}
\section*{Câu 8. Tương tác nhiều agent và tính không dừng}
\addcontentsline{toc}{section}{Câu 8. Tương tác nhiều agent và tính không dừng}
Hành động của đối thủ ảnh hưởng đến agent theo hai đường: (i) \emph{phần thưởng} --
khi đối thủ giảm giá, $p_i$ giảm 0{,}10, ví dụ từ 55 xuống 45 khi nhu cầu cao;
(ii) \emph{quan sát} -- $c_i(t+1)=a_{-i}(t)$ quyết định agent ở trạng thái nào tiếp
theo. Khi đối thủ là agent luật, quan hệ này cố định, nên với agent RL môi trường
là dừng (dù không quan sát đầy đủ). Khi \emph{cả hai cùng học}, chính sách đối thủ
thay đổi theo thời gian; cùng một cặp $(s,a)$ có phân phối thưởng và trạng thái kế
tiếp khác nhau ở các giai đoạn khác nhau. Giả định Markov dừng mà Q-learning cần để
hội tụ bị vi phạm: mỗi agent đuổi theo một mục tiêu đang dịch chuyển, và các lần
khám phá của agent này trở thành nhiễu trong dữ liệu của agent kia.

Chương trình kiểm tra điều này bằng cách cho hai agent Q-learning độc lập cùng học
(10 seed, cùng tham số baseline), rồi so sánh với RL học trước đối thủ cố định:

\begin{table}[H]
\centering
\small
\caption{Đối thủ cố định và hai agent cùng học (đánh giá tham lam, 10 seed).}
\label{tab:marl}
\begin{tabular}{lrrr}
\toprule
Thiết lập & Thưởng ĐG (TB $\pm$ SD) & Tỷ lệ giảm giá & Đổi CS (nửa sau) \\
\midrule
RL vs luật cố định (ghế RL) & $34{,}724\pm0{,}747$ & 0{,}040 & 256{,}3 \\
RL vs RL (20 agent) & $33{,}266\pm2{,}469$ & 0{,}187 & 284{,}4 \\
\bottomrule
\end{tabular}
\end{table}

\begin{Output}
seed  11 | policies 0100 / 0000 | eval = 35.19 / 35.26 | changes = 271 / 271
seed  44 | policies 1010 / 1101 | eval = 28.73 / 29.90 | changes = 272 / 271
seed  66 | policies 0000 / 1000 | eval = 32.27 / 34.20 | changes = 288 / 297
seed  77 | policies 1100 / 0100 | eval = 33.42 / 32.64 | changes = 247 / 262
seed  99 | policies 0011 / 0100 | eval = 27.31 / 29.62 | changes = 294 / 291
seed 110 | policies 0100 / 1100 | eval = 32.89 / 33.76 | changes = 275 / 296
\end{Output}

\Fig{simultaneous_learning.png}{Câu 8: thưởng và tỷ lệ giảm giá trong huấn luyện khi đối thủ cố định và khi hai agent cùng học (TB 10 seed).}

Khi cùng học, 5/10 seed kết thúc ở cặp chính sách kém (seed 44, 66, 77, 99, 110),
so với 1/10 khi đối thủ cố định. Lợi nhuận đánh giá giảm 1{,}46 mỗi bước, độ lệch
chuẩn tăng hơn ba lần, và tỷ lệ giảm giá trong huấn luyện luôn cao hơn
(0{,}20--0{,}30 so với 0{,}10--0{,}20). Ở seed 44 và 99 cả hai agent cùng giảm giá ở
nhiều trạng thái, một dạng ``chiến tranh giá'' nhỏ làm cả hai mất lợi nhuận. Số lần
đổi chính sách chỉ tăng nhẹ (284 so với 256); phần lớn dao động đã có sẵn do nhiễu
(Thí nghiệm B), nên thí nghiệm này chưa tách riêng được phần do tính không dừng.

\textbf{Đề xuất thí nghiệm để nghiên cứu sâu hơn:} \emph{học luân phiên} -- đóng
băng agent 2 trong $M$ episode trong khi agent 1 học, rồi đổi vai. Mỗi pha khi đó là
bài toán dừng. So sánh với học đồng thời (cùng tổng số bước, cùng seed) về số seed
đạt cặp chính sách tốt và tần suất đổi chính sách sẽ cho biết phần suy giảm nào do
tính không dừng. Một biến thể khác là đưa hành động trước của chính agent vào trạng
thái, hoặc cho agent học trên hành động chung (joint-action learner).

\needspace{10\baselineskip}
\section*{Câu 9. Hạn chế và độ bền vững}
\addcontentsline{toc}{section}{Câu 9. Hạn chế và độ bền vững}
\textbf{Hạn chế của mô phỏng tham chiếu:}
\begin{enumerate}
\item \textbf{Trạng thái không Markov với agent RL.} Agent luật phản ứng với hành
  động trước của RL, nhưng trạng thái $(d,c)$ của RL không chứa hành động trước
  của chính nó. Hậu quả của việc giảm giá (đối thủ giảm giá ở bước sau) bị gán cho
  một trạng thái không liên quan. Trạng thái cũng không chứa số bước còn lại trong
  episode, trong khi bước cuối dùng mục tiêu khác.
\item \textbf{Mô hình chuyển đổi tuyến tính và minh họa:} hệ số 0{,}15 và 0{,}10
  không được ước lượng từ dữ liệu; nhu cầu chỉ có hai mức và độc lập với hành động
  (giảm giá không kéo khách quay lại sau đó, không có hiệu ứng mùa vụ).
\item \textbf{Chỉ hai hành động và một mức giảm giá;} không có tồn kho, chi phí
  vận hành, khách hàng trung thành hay khác biệt sản phẩm.
\item \textbf{Đối thủ cố định và đối xứng:} RL học trước một luật duy nhất; chính
  sách học được có thể kém khi đối thủ đổi cách chơi.
\item \textbf{Kết luận nhạy với tham số:} chính sách tối ưu đổi hẳn khi mức giảm giá
  đổi từ 20\% xuống 5\% (Bảng~\ref{tab:exact}).
\end{enumerate}

\textbf{Thay đổi được cài đặt.} Hàm \texttt{scenario\_list} thay đổi \emph{mức giảm
giá} (5\%, 10\%, 30\%), \emph{hàm thưởng} (doanh thu) và \emph{hành vi đối thủ}
(tit-for-tat: lặp lại hành động trước của RL bất kể nhu cầu). Mỗi kịch bản được huấn
luyện và đánh giá trên 10 seed như Thí nghiệm E; mức tối ưu được tính chính xác.
\CodeLines{1567}{1584}

\begin{table}[H]
\centering
\small
\caption{Giá trị kỳ vọng chính xác mỗi bước của các chính sách mốc trong từng kịch bản.}
\label{tab:exact}
\begin{tabular}{lrrrrc}
\toprule
Kịch bản & Không giảm & Luôn giảm & Luật & Tối ưu $(d,c)$ & Chính sách tối ưu \\
\midrule
Giảm 20\% (tham chiếu) & 35{,}000 & 24{,}900 & 35{,}000 & 35{,}000 & 0000 \\
Giảm 5\% & 35{,}000 & 37{,}350 & 35{,}000 & 37{,}350 & 1111 \\
Giảm 10\% & 35{,}000 & 33{,}200 & 35{,}000 & 35{,}000 & 0000 \\
Giảm 30\% & 35{,}000 & 16{,}600 & 35{,}000 & 35{,}000 & 0000 \\
Thưởng doanh thu, 20\% & 70{,}000 & 66{,}400 & 70{,}000 & 70{,}000 & 0000 \\
Giảm 5\%, đối thủ tit-for-tat & 35{,}000 & 35{,}212 & 35{,}000 & 35{,}675 & 0011 \\
\bottomrule
\end{tabular}
\end{table}

\begin{table}[H]
\centering
\small
\caption{Lợi nhuận đánh giá mỗi bước (TB $\pm$ SD, 10 seed) và số seed đạt tối ưu. Kịch bản doanh thu báo lợi nhuận thực tế, không phải doanh thu.}
\label{tab:scenario}
\begin{tabular}{lcccc}
\toprule
Kịch bản & Luật & RL $(d,c)$ & RL $(d,c,\text{own})$ & Dựa trên mô hình \\
\midrule
Giảm 20\% & $34{,}99\pm0{,}18$ & $34{,}72\pm0{,}75$ (9) & $33{,}07\pm1{,}91$ (4) & $34{,}99\pm0{,}18$ (10) \\
Giảm 5\% & $34{,}99\pm0{,}18$ & $36{,}72\pm0{,}70$ (2) & $36{,}51\pm0{,}62$ (1) & $37{,}33\pm0{,}14$ (9) \\
Giảm 10\% & $34{,}99\pm0{,}18$ & $33{,}85\pm0{,}59$ (1) & $33{,}87\pm0{,}39$ (1) & $34{,}99\pm0{,}18$ (10) \\
Giảm 30\% & $34{,}99\pm0{,}18$ & $34{,}54\pm1{,}32$ (9) & $34{,}03\pm2{,}20$ (8) & $34{,}99\pm0{,}18$ (10) \\
Doanh thu, 20\% & $34{,}99\pm0{,}18$ & $29{,}07\pm3{,}20$ (1) & $29{,}79\pm2{,}24$ (1) & $34{,}99\pm0{,}18$ (10) \\
5\%, tit-for-tat & $34{,}99\pm0{,}18$ & $35{,}16\pm0{,}46$ (1) & $34{,}98\pm0{,}38$ (0) & $35{,}59\pm0{,}18$ (5) \\
\bottomrule
\end{tabular}
\end{table}

\Fig{scenario_comparison.png}{Câu 9--10: lợi nhuận đánh giá khi thị trường thay đổi.\label{fig:scenario}}

\textbf{Ảnh hưởng của thay đổi.}
\begin{itemize}
\item \textbf{Giảm 5\%:} lợi nhuận mỗi đơn khi giảm giá là 4{,}5. Khi nhu cầu thấp,
  giảm giá cho $10\times0{,}45\times4{,}5=20{,}25>15$; khi nhu cầu cao,
  $20\times0{,}70\times4{,}5=63>55$. Luôn giảm giá trở thành tối ưu (37{,}35). Agent
  luật không đổi nên mất 2{,}35 mỗi bước (6{,}3\%); RL tham chiếu tự chuyển sang giảm
  giá 75\% thời gian và đạt 36{,}72. Đây là kịch bản cho thấy rõ lợi thế thích ứng
  của RL.
\item \textbf{Giảm 10\%:} giảm giá có lợi trước mắt (khi nhu cầu thấp 18 so với 15;
  khi nhu cầu cao 56 so với 55) nhưng kích hoạt phản ứng giảm giá của đối thủ ở bước
  sau. Tối ưu vẫn là không giảm giá. RL tham chiếu chỉ thấy lợi ích trước mắt, vì hậu
  quả xảy ra ở trạng thái mà nó không phân biệt được (hạn chế 1). Kết quả là RL giảm
  giá 50\% thời gian và đạt 33{,}85 -- \emph{thua} agent luật.
\item \textbf{Giảm 30\%:} giảm giá thiệt hại rõ (biên lợi nhuận chỉ còn 2), nên RL
  học đúng 9/10 seed, gần giống kịch bản tham chiếu.
\item \textbf{Đối thủ tit-for-tat, 5\%:} tối ưu là chỉ giảm giá khi nhu cầu cao
  (35{,}675). RL tham chiếu đạt 35{,}16, nhỉnh hơn luật nhưng chỉ 1/10 seed đạt tối
  ưu, vì khoảng cách giá trị giữa các chính sách rất nhỏ so với nhiễu.
\end{itemize}

\needspace{10\baselineskip}
\section*{Câu 10. Thiết kế agent tốt hơn}
\addcontentsline{toc}{section}{Câu 10. Thiết kế agent tốt hơn}
Phân tích ở câu 7--9 chỉ ra hai điểm yếu của agent RL tham chiếu: \emph{trạng thái
thiếu thông tin} và \emph{cập nhật TD với $\alpha$ cố định quá nhiễu}. Agent đề xuất
(\texttt{ModelBasedAgent}) xử lý cả hai.

\begin{itemize}
\item \textbf{Trạng thái:} $s_i=(d_t,\,a_{-i}(t-1),\,a_i(t-1))$, 8 trạng thái. Thêm
  hành động trước của chính agent làm trạng thái trở thành Markov khi đối thủ phản
  ứng theo hành động trước (agent luật, tit-for-tat).
\item \textbf{Hành động:} $\{0,1\}$ như đề bài; có thể mở rộng thành nhiều mức giảm
  giá mà không đổi thuật toán.
\item \textbf{Hàm thưởng:} lợi nhuận~\eqref{eq:reward}, đúng mục tiêu kinh doanh
  (câu 5 cho thấy thưởng doanh thu dẫn tới hành vi sai).
\item \textbf{Phương pháp học và ra quyết định:} học \emph{mô hình} thay vì học trực
  tiếp $Q$. Agent đếm $N(s,a)$, tổng thưởng và tần suất trạng thái kế tiếp để ước
  lượng $\hat R(s,a)$ và $\hat P(s'\mid s,a)$. $\hat P$ chứa cả quá trình nhu cầu lẫn
  \emph{mô hình phản ứng của đối thủ}. Cứ 50 episode và khi kết thúc huấn luyện,
  agent giải $Q(s,a)=\hat R(s,a)+\gamma\sum_{s'}\hat P(s'\mid s,a)\max_{a'}Q(s',a')$
  bằng lặp giá trị với $\gamma=0{,}98$, gần mục tiêu thưởng trung bình hơn
  $\gamma=0{,}9$. Ước lượng bằng trung bình mẫu có phương sai giảm dần theo $1/N$,
  thay vì giữ mức nhiễu cố định như $\alpha=0{,}1$. Khám phá dùng cùng lịch
  $\varepsilon$ với baseline để so sánh công bằng.
\item \textbf{Quy trình đánh giá:} cùng giao thức tách biệt như Thí nghiệm E (10
  seed, 500 episode, $\varepsilon=0$, số ngẫu nhiên chung), cộng giá trị kỳ vọng
  chính xác và mức tối ưu cho từng kịch bản. Để đo độ bền vững, dùng
  \emph{regret} = tối ưu $-$ giá trị kỳ vọng, và lấy \emph{regret lớn nhất} trên tập
  kịch bản đa dạng (mức giảm giá, hàm thưởng, đối thủ). Agent bền vững hơn là agent
  có regret lớn nhất nhỏ hơn và độ lệch chuẩn giữa các seed nhỏ hơn.
\end{itemize}

\CodeLines{342}{403}

\begin{table}[H]
\centering
\small
\caption{Regret = tối ưu $-$ giá trị kỳ vọng chính xác (TB 10 seed). Kịch bản doanh thu tính theo đơn vị doanh thu.}
\label{tab:regret}
\begin{tabular}{lcccc}
\toprule
Kịch bản & Luật & RL $(d,c)$ & RL $(d,c,\text{own})$ & Dựa trên mô hình \\
\midrule
Giảm 20\% & 0{,}000 & 0{,}255 & 1{,}887 & 0{,}000 \\
Giảm 5\% & 2{,}350 & 0{,}630 & 0{,}868 & 0{,}020 \\
Giảm 10\% & 0{,}000 & 1{,}149 & 1{,}121 & 0{,}000 \\
Giảm 30\% & 0{,}000 & 0{,}439 & 0{,}952 & 0{,}000 \\
Doanh thu, 20\% & 0{,}000 & 2{,}298 & 2{,}242 & 0{,}000 \\
5\%, tit-for-tat & 0{,}675 & 0{,}539 & 0{,}717 & 0{,}085 \\
\midrule
\textbf{Regret lớn nhất} & 2{,}350 & 2{,}298 & 2{,}242 & \textbf{0{,}085} \\
\bottomrule
\end{tabular}
\end{table}

\textbf{Kết quả kiểm tra.} Agent dựa trên mô hình đạt tối ưu ở 10/10 seed trong bốn
kịch bản, 9/10 ở kịch bản 5\% (trung bình 37{,}33 so với tối ưu 37{,}35) và 5/10 với
đối thủ tit-for-tat (35{,}59 so với 35{,}675). Regret lớn nhất chỉ 0{,}085, so với
2{,}35 của luật (thua khi giảm 5\%) và khoảng 2{,}3 của RL tham chiếu (thua khi khoảng
cách giá trị nhỏ). Chính sách của nó gần như không đổi trong nửa sau huấn luyện
(0--2{,}3 lần, so với hàng trăm lần của Q-learning).

Một kết quả đối chứng quan trọng: \emph{chỉ thêm trạng thái mà vẫn dùng Q-learning}
(\texttt{ExtendedQLearningAgent}) không cải thiện, thậm chí kém hơn ở kịch bản
tham chiếu (33{,}07, 4/10). Bảng Q có 16 giá trị nên mỗi ô nhận ít dữ liệu hơn, trong
khi nhiễu cập nhật không giảm. Như vậy, cải thiện của agent đề xuất đến từ việc
\emph{kết hợp} trạng thái đầy đủ với ước lượng ít nhiễu và lập kế hoạch, không đến
từ riêng phần trạng thái.

\textbf{Giới hạn của đề xuất:} agent cần không gian trạng thái nhỏ, rời rạc để
đếm; nó giả định đối thủ dừng trong lúc học (với đối thủ cũng đang học, mô hình cần
quên dữ liệu cũ, ví dụ dùng cửa sổ trượt); và tập kịch bản kiểm tra ở đây vẫn nhỏ.
Một kiểm tra mạnh hơn sẽ thêm đối thủ ngẫu nhiên, đối thủ đổi chiến lược giữa chừng
và nhu cầu có tính mùa vụ, rồi so sánh regret lớn nhất của cả ba agent trên tập đó.

% ============================================================
\chapter{Thảo luận, hạn chế và đề xuất}
\label{chap:discussion}

\section{Phân biệt quan sát và giải thích}
\begin{table}[H]
\centering
\small
\caption{Quan sát đo được và mức độ chắc chắn của giải thích.}
\begin{tabularx}{\textwidth}{XX>{\raggedright\arraybackslash}p{4.3cm}}
\toprule
Quan sát & Giải thích & Mức bằng chứng \\
\midrule
RL hòa luật ở 9/10 seed (tham chiếu) & Luật trùng chính sách tối ưu; RL tốt nhất chỉ hòa & Chứng minh bằng liệt kê chính xác \\
Thưởng huấn luyện không tăng sau $\sim$500 episode & Chính sách tham lam dao động giữa tối ưu và kém & Đo trực tiếp (Hình~\ref{fig:pv}) \\
Dao động lớn của $Q$ & Nhiễu nhị thức, $\alpha$ cố định, trạng thái thiếu số bước còn lại & Hợp lý, chưa tách riêng từng nguyên nhân \\
$\varepsilon=0$ khóa vào lựa chọn ban đầu & Thưởng luôn dương, $Q$ khởi tạo 0 & Suy ra trực tiếp từ công thức cập nhật \\
$\varepsilon=0{,}30$ cho chính sách tham lam kém & Điểm hội tụ phụ thuộc phân phối hành vi khi trạng thái không Markov & Giả thuyết, cần thí nghiệm riêng \\
RL thua luật khi giảm 10\% & RL không thấy hậu quả trễ của giảm giá & Phù hợp với tính toán; agent dựa trên mô hình sửa được \\
Hai agent cùng học kém hơn & Không dừng và vòng phản hồi giảm giá & Quan sát; chưa tách khỏi nhiễu \\
\bottomrule
\end{tabularx}
\end{table}

\section{Các nhận xét chung}
\begin{itemize}
\item \textbf{Agent luật} cho quyết định nhất quán (độ lệch chuẩn giữa các seed chỉ
  0{,}18, hoàn toàn do nhiễu doanh số), nhưng hiệu quả phụ thuộc vào việc luật có
  khớp với môi trường hay không: tối ưu ở mức giảm 10--30\%, mất 6{,}3\% lợi nhuận
  ở mức 5\%.
\item \textbf{Agent RL} cải thiện được qua phản hồi và thích ứng khi thị trường đổi,
  nhưng việc học không đảm bảo chính sách tối ưu: kết quả phụ thuộc seed,
  $\alpha$, $\varepsilon$ và cách biểu diễn trạng thái.
\item \textbf{Nhiễu ngẫu nhiên} làm thưởng dao động ngay cả khi chính sách không đổi
  (34{,}77--35{,}26 giữa các seed đánh giá của cùng chính sách). Vì vậy giá trị kỳ
  vọng chính xác và số ngẫu nhiên chung rất hữu ích để so sánh.
\item \textbf{Hai agent cùng học} ảnh hưởng đến kết quả của nhau (câu 8).
\item \textbf{Huấn luyện và đánh giá đo hai thứ khác nhau:} thưởng huấn luyện gồm
  chi phí khám phá và chính sách đang thay đổi; thưởng đánh giá đo chính sách cuối
  với bảng Q đóng băng.
\end{itemize}

\section{Hạn chế của chính báo cáo này}
\begin{itemize}
\item 10 seed đủ để thấy xu hướng nhưng khoảng tin cậy vẫn rộng; các so sánh như
  $\alpha=0{,}01$ với $0{,}10$ không tách biệt được về mặt thống kê.
\item Agent đề xuất được thiết kế \emph{sau khi} đã thấy điểm yếu của agent tham
  chiếu trên chính các kịch bản này, nên có nguy cơ thiết kế khớp quá mức với tập
  kiểm tra. Cần thêm kịch bản mới chưa dùng khi thiết kế để kiểm tra lại.
\item Mọi kết luận chỉ đúng trong mô hình mô phỏng đơn giản hóa; mô hình chuyển đổi
  là minh họa, không được kiểm chứng với dữ liệu thật \cite{complementary}.
\end{itemize}

\section{Đề xuất cải tiến}
\begin{enumerate}
\item Đưa hành động trước của chính agent (và nếu cần, số bước còn lại) vào trạng
  thái; dùng $\alpha$ giảm theo số lần thăm hoặc học dựa trên mô hình.
\item Mở rộng môi trường: nhu cầu có quán tính (chuỗi Markov), nhiều mức giảm giá,
  tồn kho, khách hàng trung thành để giảm giá có lợi ích dài hạn.
\item Đánh giá theo regret trên một tập kịch bản và đối thủ đa dạng, báo cáo trung
  bình, độ lệch chuẩn và trường hợp xấu nhất.
\item Với nhiều agent cùng học: học luân phiên, mô hình hóa đối thủ có quên dữ liệu
  cũ, hoặc các thuật toán học đa agent như joint-action learner.
\end{enumerate}

% ============================================================
\chapter{Kết luận}
\label{chap:conclusion}
Báo cáo đã cài đặt môi trường thương mại điện tử hai người bán, agent Q-learning và
agent dựa trên luật, tách riêng giai đoạn huấn luyện và đánh giá, và thực hiện đủ
năm thí nghiệm A--E trên 10 seed độc lập cùng ba phần mở rộng cho câu 8--10.

Trong môi trường tham chiếu, kết quả đo được không ủng hộ nhận định RL tốt hơn hay kém
hơn luật: luật trùng với chính sách tối ưu (35{,}00), RL đạt tối ưu ở 9/10 seed và
chênh lệch không có ý nghĩa thống kê. Kết quả đổi khi môi trường đổi: ở mức giảm giá
5\%, RL thích ứng và vượt luật; ở mức 10\%, RL thua luật vì trạng thái không thể hiện
hậu quả trễ của việc giảm giá. Tốc độ học và lịch khám phá ảnh hưởng mạnh đến độ ổn
định: $\alpha=0{,}5$ và $\varepsilon=0{,}30$ cố định cho chính sách kém hơn nhiều;
$\varepsilon=0$ khóa agent vào lựa chọn ngẫu nhiên ban đầu. Khi hai agent cùng học,
số seed kết thúc ở chính sách kém tăng từ 1/10 lên 5/10.

Agent đề xuất ở câu 10 -- trạng thái có hành động trước của chính agent, học mô hình
thị trường và đối thủ bằng đếm, rồi lập kế hoạch -- có regret lớn nhất 0{,}085 trên
sáu kịch bản, so với khoảng 2{,}3 của cả agent luật lẫn agent RL tham chiếu. Bài học
chính của thí nghiệm là: chọn agent nào không thể dựa vào một lần chạy với phần
thưởng cao nhất; cần một quy trình đánh giá tách biệt, có mốc tối ưu, lặp nhiều seed
và thử trên nhiều điều kiện môi trường.

\clearpage
\begin{thebibliography}{9}
\bibitem{assignment} Intelligent System Development, \emph{Laboratory Experiment 02:
Reinforcement-Learning-Based Agents versus Rule-Based Agents}, năm học 2026--2027,
tệp \path{TAILIEU/ASSIGNMENT-AGENT_02/assignment_Agent02_due_29.10.pdf}.
\bibitem{complementary} Intelligent System Development, \emph{AGENT02: Python
Implementation -- Reinforcement-Learning-Based Agents and Rule-Based Agents},
tệp \path{TAILIEU/ASSIGNMENT-AGENT_02/agent02_complementary_code.pdf} và mã tham khảo
\path{agent02_experiment.py}.
\bibitem{sutton} R. S. Sutton và A. G. Barto, \emph{Reinforcement Learning: An
Introduction}, 2nd ed., MIT Press, 2018.
\bibitem{code} Phan Văn Khôi, \emph{agent02\_experiment.py} và kết quả thí nghiệm,
\url{https://github.com/KhoiF/PTHTTM}, thư mục \path{src/Agent 02/}.
\end{thebibliography}
\end{document}
